% ============================================================================
% Fichier  : partie4/P04_C18_forensique_iot_embarque.tex
% Titre    : Forensique IoT et Systèmes Embarqués
% Auteur   : MINKA MI NGUIDJOÏ Thierry Emmanuel
% Version  : 1.0 -- Juin 2026
% Statut   : RÉDIGÉ
% Sources  : Construit ex-nihilo (aucune source projet ne couvre l'IoT)
%            Ancré dans le contexte africain/camerounais : routeurs ADSL,
%            caméras IP, Box opérateurs (MTN/Orange), compteurs prépayés
%            Référence : NIST SP 800-213 (IoT Security), IEEE P2933,
%            ENISA IoT Forensics Guidelines
% Positionnement : Chapitre le plus avancé de la Partie IV. Niveau Expert
%                  justifié : UART/JTAG, extraction firmware, analyse binaire.
%                  Contexte camerounais fort : cyberattaques sur routeurs
%                  opérateurs, caméras SONATREL/aéroport, Box DSTV/Canal+.
% ============================================================================

\chapter{Forensique IoT et Systèmes Embarqués}
\label{chap:for-iot}

\epigraph{%
    « The Internet of Things is the Internet of Vulnerabilities. »%
}{--- Bruce Schneier}

\niveauChapitre{Expert}{%
    Notions d'électronique embarquée, Linux embarqué et protocoles
    série (UART/I2C/SPI) souhaitables. Ce chapitre est le plus
    technique de la Partie~IV.%
}

\begin{boxobjectifs}
\begin{itemize}
    \item Comprendre les architectures IoT typiques et identifier
          où résident les données forensiquement utiles.
    \item Extraire les logs d'un routeur compromis via son interface
          d'administration ou via UART.
    \item Analyser le firmware d'un appareil IoT~: identifier les
          systèmes de fichiers, les configurations, les backdoors.
    \item Investiguer les caméras IP et systèmes DVR/NVR~:
          extraction des enregistrements et logs d'accès.
    \item Appliquer les principes de custody et d'intégrité à des
          preuves IoT~: volatilité extrême, absence de write blocker
          standardisé.
    \item Identifier les spécificités IoT du contexte camerounais~:
          Box opérateurs, compteurs prépayés, caméras de surveillance.
\end{itemize}
\end{boxobjectifs}

\bigskip

% ============================================================================
\section*{L'IoT au Cameroun~: contexte forensique}
% ============================================================================

L'IoT forensique n'est pas encore au c\oe ur des investigations camerounaises,
mais plusieurs incidents récents imposent de s'y préparer~:

\begin{itemize}
    \item Routeurs ADSL/4G (MTN Box, Orange HomeBox) compromis pour
          des attaques Man-in-the-Middle sur les clients~;
    \item Caméras IP de surveillance (hôtels, banques, aéroport
          de Nsimalen) piratées~; enregistrements utilisés pour
          du chantage ou effacés après un incident~;
    \item Compteurs électriques prépayés SONATREL compromis
          par manipulation des compteurs~;
    \item Décodeurs Canal+/DSTV modifiés pour contourner les
          protections DRM~: dommages aux droits d'auteur.
\end{itemize}

La particularité forensique de ces appareils~: ils produisent des
preuves importantes (logs de connexion, enregistrements vidéo,
configuration réseau) mais leur acquisition requiert des techniques
très différentes de celles couvertes dans les chapitres précédents.

\separateur

% ============================================================================
\section{Architecture IoT~: où Résident les Données Forensiques}
\label{sec:architecture-iot}
% ============================================================================

\subsection{Composants typiques d'un appareil IoT}

\begin{tableaucours}{p{2.8cm}p{3.5cm}p{5.5cm}}{%
    \tableauHeader{Composant} &
    \tableauHeader{Type de mémoire} &
    \tableauHeader{Données forensiques}
}
\textbf{Flash NAND/NOR} & Non volatile (persiste hors tension) &
    Firmware, configuration, logs persistants, données utilisateur \\
\textbf{RAM (DRAM/SRAM)} & Volatile (perdue à l'extinction) &
    Processus actifs, connexions réseau, clés de session,
    configuration temporaire \\
\textbf{eMMC} & Non volatile (stockage embarqué) &
    Système d'exploitation Linux, journaux, historique d'activité \\
\textbf{SD Card} & Non volatile (amovible) &
    Enregistrements (caméras IP), logs (DVR), sauvegarde config \\
\textbf{EEPROM} & Non volatile & Configuration critique, clés, calibrations \\
\end{tableaucours}
\vspace{-0.8em}
\begin{center}{\small\itshape Tableau~18.1 --- Mémoires IoT et données forensiques}\end{center}

\subsection{Priorité de collecte~: RFC~3227 appliquée à l'IoT}

La RFC~3227 (ordre de volatilité, Ch.~\ref{chap:meth-std}) s'applique
à l'IoT avec des adaptations~:

\begin{enumerate}
    \item \textbf{RAM et processus actifs}~: capturer via interface
          d'administration (SSH/Telnet/API) \emph{avant} toute autre action.
          Durée de vie~: de quelques secondes (si l'appareil est en train
          d'être réinitialisé) à des jours.
    \item \textbf{Logs système en mémoire}~: \texttt{dmesg},
          \texttt{syslog} en RAM. Exportés via interface admin avant
          extinction.
    \item \textbf{Logs persistants sur flash}~: survivent à l'extinction,
          mais souvent effacés au redémarrage ou en rotation courte.
    \item \textbf{Configuration}~: paramètres réseau, identifiants,
          règles de pare-feu. Non volatile mais peut être modifiée.
    \item \textbf{Firmware complet}~: extraction via UART ou déchippage~;
          analyse offline.
\end{enumerate}

\begin{boxalert}
\textbf{Le piège du ``factory reset''}

La première réaction d'un suspect en possession d'un appareil IoT
compromis est souvent d'effectuer un reset d'usine. Contrairement
aux idées reçues~:
\begin{itemize}
    \item Le factory reset sur la plupart des routeurs efface les logs
          et la configuration en RAM, mais \textbf{pas nécessairement
          le firmware} ni les données de la partition système~;
    \item Sur les caméras IP avec SD~Card, la SD~Card n'est pas
          systématiquement effacée par le reset (elle est montée
          séparément)~;
    \item L'enregistrement du reset lui-même peut apparaître dans les
          logs de l'opérateur (réinitialisation du bail DHCP,
          reconnexion réseau).
\end{itemize}
\textbf{Protocole~:} saisir physiquement l'appareil \emph{avant}
toute tentative de reset~; retirer la carte SD si présente~; couper
l'alimentation proprement (pas d'arrêt brutal qui pourrait corrompre
les données flash).
\end{boxalert}

% ============================================================================
\section{Forensique Routeur~: la Cible Prioritaire}
\label{sec:forensique-routeur}
% ============================================================================

\subsection{Pourquoi les routeurs sont des cibles forensiques clés}

Un routeur compromis est idéalement placé pour~:
\begin{itemize}
    \item Voir \textbf{tout le trafic} passant par lui (Man-in-the-Middle)~;
    \item Rediriger le DNS pour du phishing (DNS poisoning)~;
    \item Exfiltrer les credentials Wi-Fi de tous les clients~;
    \item Servir de rebond pour des attaques vers d'autres cibles
          (point de pivot).
\end{itemize}

Un routeur \emph{victime} d'une attaque contient les logs qui prouvent
l'intrusion. Un routeur \emph{appartenant au suspect} contient les
logs de ses communications.

\subsection{Extraction via interface d'administration}

\begin{boxoutils}{Extraction de logs via interface admin (SSH)}
\begin{lstlisting}[language=bash_forensique]
# Connexion SSH au routeur (si SSH est actif)
# PREREQUIS : credentials connus + autorisation judiciaire art. 52

ssh admin@192.168.1.1

# Routeurs sous OpenWRT / DD-WRT (Linux embarque)
# Lister les processus actifs
cat /proc/running_processes
# Ou
ps aux

# Connexions reseau actives (equivalent netstat)
cat /proc/net/tcp   # connexions TCP
cat /proc/net/arp   # table ARP (machines connectees recemment)
cat /proc/net/dev   # interfaces reseau et statistiques

# Logs systeme
cat /tmp/syslog    # syslog en RAM (volatile !)
logread            # commande OpenWRT pour les logs systeme

# Configuration actuelle (non volatile)
cat /etc/config/network    # configuration reseau
cat /etc/config/firewall   # regles de pare-feu
cat /etc/config/dhcp       # baux DHCP (machines ayant eu une IP)
cat /etc/config/wireless   # configuration Wi-Fi (SSID, mot de passe)

# Historique DHCP (clients ayant ete connectes)
cat /tmp/dhcp.leases
# Format : timestamp MAC-address IP hostname
# Exemple :
# 1710288900 aa:bb:cc:dd:ee:ff 192.168.1.42 LAPTOP-SUSPECT
# 1710289200 11:22:33:44:55:66 192.168.1.55 iPhone-MBONGO

# EXPORTER IMMEDIATEMENT avant toute action
scp admin@192.168.1.1:/tmp/syslog /output/router_syslog.txt
scp admin@192.168.1.1:/tmp/dhcp.leases /output/router_dhcp.txt
\end{lstlisting}
\end{boxoutils}

\subsection{Extraction via UART~: quand l'interface admin est inaccessible}

Quand les credentials de l'interface admin sont inconnus (changés par
l'attaquant, ou appareil saisi dont le propriétaire refuse de coopérer),
l'accès physique via \termecle{UART}\index{UART} permet d'obtenir une
console série directe sur l'appareil.

\begin{boxdef}
\textbf{UART} (\textit{Universal Asynchronous Receiver-Transmitter})~:
interface série physique présente sur la quasi-totalité des cartes
électroniques embarquées. Elle donne accès à une console root lors
du démarrage du bootloader (U-Boot) et souvent tout au long du
fonctionnement. Fonctionne via 3 broches~: TX (émission), RX
(réception), GND (masse).
\end{boxdef}

\begin{boxoutils}{Accès UART à un routeur}
\begin{lstlisting}[language=bash_forensique]
# Matériel nécessaire :
# - Convertisseur USB-to-UART (CH340G ou FT232RL, ~5 USD)
# - Voltmetre (pour mesurer le niveau logique : 3.3V ou 5V)
# - Fils Dupont / sonde oscilloscope

# Etape 1 : Ouvrir le boitier et localiser les broches UART
# Sur la plupart des routeurs : groupe de 4 broches (VCC, TX, RX, GND)
# Identifier GND (masse) avec le voltmetre

# Etape 2 : Mesurer le niveau logique
# Si VCC mesure 3.3V -> utiliser un adaptateur 3.3V
# Si VCC mesure 5V -> adaptateur 5V ou niveau-shifter 5V->3.3V

# Etape 3 : Connecter (CROISER TX<->RX)
# Convertisseur TX -> Routeur RX
# Convertisseur RX -> Routeur TX
# Convertisseur GND -> Routeur GND

# Etape 4 : Terminal serie (Linux)
screen /dev/ttyUSB0 115200    # vitesse la plus courante
# Alternatives : minicom, picocom
picocom -b 115200 /dev/ttyUSB0

# Resultat : console U-Boot puis Linux au demarrage
# [    0.000000] Linux version 5.4.152 (OpenWrt)
# [    2.341827] eth0: Link is Up - 1Gbps/Full
# OpenWrt login: root     # <- shell root accessible !

# Commande de sauvegarde complete depuis UART
# (lente mais complete : dump de la flash NAND entiere)
cat /dev/mtdblock0 > /tmp/firmware_dump.bin
# Puis exfiltrer via SCP ou TFTP
tftp -p -l /tmp/firmware_dump.bin 192.168.1.100
\end{lstlisting}
\end{boxoutils}

\begin{boiteRemarque}{UART et responsabilité}
L'accès UART à un appareil implique son ouverture physique (vis, parfois
clips) et la manipulation de son électronique. Documenter photographiquement
l'état de l'appareil \emph{avant} toute manipulation~: position des
broches, état de la carte, numéro de série visible. Toute modification
physique doit être mentionnée dans le rapport (ACPO Principe~2~: si une
modification est nécessaire, elle doit être documentée et justifiée).
\end{boiteRemarque}

% ============================================================================
\section{Analyse de Firmware~: ce que l'Appareil Cache}
\label{sec:analyse-firmware}
% ============================================================================

\subsection{Extraction et analyse avec Binwalk}

\termecle{Binwalk}\index{Binwalk} est l'outil de référence pour l'analyse
de firmware. Il identifie et extrait automatiquement les systèmes de
fichiers, les images compressées et les signatures dans un binaire.

\begin{boxoutils}{Analyse de firmware avec Binwalk}
\begin{lstlisting}[language=bash_forensique]
# Analyser un fichier firmware (telecharge sur le site du fabricant
# ou extrait de l'appareil)
binwalk firmware.bin

# Resultat type :
# DECIMAL   HEX       DESCRIPTION
# 0         0x0       TRX firmware header
# 28        0x1C      LZMA compressed data
# 131072    0x20000   Squashfs filesystem, little endian, version 4.0
#                     Compression: xz
#                     Total inodes: 2847
#                     Block size: 262144

# Extraire automatiquement tous les composants identifie
binwalk -e --directory=/output/firmware_extracted/ firmware.bin

# Contenu type apres extraction :
ls /output/firmware_extracted/
# _firmware.bin.extracted/
#   0.bin          <- bootloader
#   20000.squashfs <- systeme de fichiers (monter pour analyser)

# Monter le systeme de fichiers SquashFS
sudo mount -o loop,ro \
    /output/firmware_extracted/_firmware.bin.extracted/20000.squashfs \
    /mnt/firmware/

# Explorer le contenu
ls /mnt/firmware/
# bin/  etc/  lib/  overlay/  proc/  sbin/  usr/  var/

# Chercher les configurations et identifiants
cat /mnt/firmware/etc/shadow      # hash des mots de passe
cat /mnt/firmware/etc/passwd      # comptes utilisateurs
cat /mnt/firmware/etc/config/     # configuration OpenWRT
find /mnt/firmware/ -name "*.conf" -exec grep -l "password\|passwd" {} \;

# Rechercher les backdoors (cles SSH pre-installees)
find /mnt/firmware/ -name "authorized_keys"
find /mnt/firmware/ -name "*.key" -o -name "*.pem"

# Identifier les services actifs au demarrage
cat /mnt/firmware/etc/rc.d/
cat /mnt/firmware/etc/inittab
\end{lstlisting}
\end{boxoutils}

\begin{boxscenario}{Backdoor découverte dans un routeur MTN Box}
L'analyse Binwalk du firmware d'un routeur MTN Box saisi dans une
investigation de cybercriminalité révèle~:

\begin{lstlisting}[language=bash_forensique]
find /mnt/firmware/ -name "authorized_keys"
# /mnt/firmware/root/.ssh/authorized_keys

cat /mnt/firmware/root/.ssh/authorized_keys
# ssh-rsa AAAAB3NzaC1yc2EAAAADAQABAAAB... attacker@remote-c2.ru
\end{lstlisting}

\textbf{Constatation~:} le fichier \texttt{/root/.ssh/authorized\_keys}
du firmware extrait contient une clé SSH publique associée à
\texttt{attacker@remote-c2.ru}. Cette clé permet un accès root sans
mot de passe au routeur depuis l'IP contrôlée par cet attaquant.\\[0.3ex]
\textbf{Hash du fichier firmware~:}
\hash{SHA-256}{a3f7b29c8e1d4f56a2c8b7e3d9f0a142b5c8d3e6f1a4b7c0}\\[0.3ex]
\textbf{Corrélation~:} l'adresse \texttt{remote-c2.ru} apparaît dans
les logs SIEM de l'entreprise victime sur les 3 derniers mois~--- le
routeur servait de backdoor persistante.
\end{boxscenario}

% ============================================================================
\section{Forensique Caméras IP et Systèmes DVR/NVR}
\label{sec:cameras-ip}
% ============================================================================

\subsection{Architecture d'un système de vidéosurveillance}

Un système de vidéosurveillance IP typique comprend~:
\begin{itemize}
    \item Des \textbf{caméras IP} (ou caméras CCTV analogiques
          convertis par des encodeurs)~;
    \item Un \textbf{NVR} (\textit{Network Video Recorder}) ou
          \textbf{DVR} (\textit{Digital Video Recorder}) qui stocke
          les enregistrements~;
    \item Un \textbf{HDD interne} (3,5 pouces SATA) dans le NVR~;
    \item Des \textbf{logs d'accès} au système de gestion vidéo (VMS).
\end{itemize}

\subsection{Extraction des enregistrements}

\begin{boxoutils}{Acquisition forensique d'un NVR/DVR}
\begin{lstlisting}[language=bash_forensique]
# OPTION 1 : Export via interface web du NVR (le plus simple)
# Connexion a l'interface web (IP:port)
# Playback -> Exporter la plage horaire d'interet -> Format H.264/MP4

# OPTION 2 : Extraction du disque dur interne
# Retirer le HDD du NVR (etiqueter le cable SATA et son emplacement)
# Photographier avant extraction
# Connecter via write blocker SATA sur le poste forensique

# Identifier le systeme de fichiers (souvent proprietaire)
fdisk -l /dev/sdb
# Resultat type :
# Disk /dev/sdb: 2 TiB
# Device     Boot   Start       End    Sectors  Size Type
# /dev/sdb1          2048   2099199    2097152    1G  Linux filesystem
# /dev/sdb2       2099200 3906250751 3904151552  1.8T  (inconnu/propriétaire)

# Si le systeme de fichiers est reconnu (ext4, FAT) :
mount -o ro /dev/sdb1 /mnt/nvr/
ls /mnt/nvr/record/
# Fichiers .dav, .mp4, .avi selon le fabricant

# Si le systeme de fichiers est proprietaire (Hikvision, Dahua)
# utiliser les outils specifiques :
# VideoInspector, DVR-Examiner, ou
# Outil Python : python-dvr (open source)
python3 dvr_examiner.py /dev/sdb2 --export /output/nvr_records/

# OPTION 3 : Carving video sur espace brut
# Si le FS est corrompu ou efface :
bulk_extractor -o /output/bulk/ /dev/sdb
# bulk_extractor detecte les en-tetes video (H.264 NAL units)
# et extrait les fragments
\end{lstlisting}
\end{boxoutils}

\subsection{Logs d'accès au VMS~: qui a visionné quoi}

\begin{boxoutils}{Analyse des logs d'accès VMS}
\begin{lstlisting}[language=bash_forensique]
# Systemes Hikvision (les plus repandus en Afrique)
# Logs dans /var/log/hikvision/ ou /mnt/mtd/log/

ssh admin@192.168.1.64
cat /mnt/mtd/log/access.log

# Contenu type :
# [2026-03-13 02:14:15] LOGIN  admin  192.168.1.1   Success
# [2026-03-13 02:14:22] PLAYBACK  admin  Camera1  2026-03-12 18:00 -> 18:30
# [2026-03-13 02:15:01] DELETE  admin  Camera1  2026-03-12 18:00 -> 18:30
# [2026-03-13 02:15:02] FORMAT  admin  HDD1

# La sequence : Login -> Playback -> Delete -> Format
# indique une tentative de destruction de preuves

# Verifier si le formatage a ete complet (les secteurs non
# réécrits peuvent contenir des donnees recuperables)
# Apres saisie du HDD :
# dd if=/dev/sdb of=/output/nvr_image.dd bs=65536
# foremost -t avi,mp4 -i /output/nvr_image.dd -o /output/nvr_carved/
\end{lstlisting}
\end{boxoutils}

\begin{boxjuris}
\textbf{Enregistrements de vidéosurveillance~: admissibilité}

En droit camerounais, les enregistrements de vidéosurveillance sont
admissibles comme preuves si~:
\begin{itemize}
    \item Le système a été installé conformément à la réglementation
          (signalétique, déclaration auprès de l'autorité compétente)~;
    \item La chaîne de custody est documentée depuis le NVR/DVR
          jusqu'à la présentation en audience~;
    \item L'intégrité des fichiers est attestée par hash SHA-256~;
    \item Un expert judiciaire atteste de l'authenticité et de
          l'horodatage (synchronisation NTP du système).
\end{itemize}
\textbf{Point de vigilance~:} vérifier la synchronisation horaire
du NVR~: une dérive de plusieurs heures invalide la corrélation
avec d'autres sources. Documenter dans le rapport~: ``Heure affichée
du NVR au moment de la saisie~: 14h22~; heure de référence GPS~:
14h25~; dérive estimée~: +3~minutes.''
\end{boxjuris}

% ============================================================================
\section{IoT Spécifiques~: Box Opérateurs et Compteurs}
\label{sec:iot-cameroun}
% ============================================================================

\subsection{Box opérateurs MTN/Orange~: forensique des données de trafic}

Les Box Internet des opérateurs camerounais (MTN HomeBox 4G,
Orange LiveBox) sont des routeurs 4G intégrant un modem,
un switch Wi-Fi et souvent un décodeur TV.

\begin{boxoutils}{Extraction des données forensiques d'une Box 4G}
\begin{lstlisting}[language=bash_forensique]
# Interface d'administration web (typiquement 192.168.1.1)
# Identifiants par defaut (si non changes) :
# MTN HomeBox : admin/admin ou admin/[IMEI partiel]
# Orange LiveBox : admin/[code sur etiquette]

# Via SSH (si active) :
ssh admin@192.168.1.1

# Historique des connexions clients Wi-Fi
cat /tmp/hostapd_events.log
# Ou via commande propriétaire :
wl assoclist              # Clients actuellement connectes
cat /proc/net/arp         # Cache ARP (connexions recentes)

# Logs DNS (si dnsmasq actif avec logging)
cat /tmp/dnsmasq.log
# Format : timestamp IP-client query type domain
# 2026-01-15 14:22:05 192.168.1.42 A medequip-portugal.com
# -> le client .42 a resolu ce domaine (phishing !)

# Volume de trafic par client (utile pour detecter exfiltration)
cat /proc/net/ip_conntrack | awk '{print $5}' | \
    cut -d= -f2 | sort | uniq -c | sort -rn | head -10

# IMPORTANT : les Box operateurs sauvegardent aussi les logs
# sur leurs serveurs centraux. Complement indispensable via
# requisition art. 57 Loi 2010/012 aupres de MTN/Orange.
\end{lstlisting}
\end{boxoutils}

\subsection{Compteurs prépayés~: manipulation et traces}

Les compteurs électriques prépayés (système SONATREL) utilisent
des tokens prépayés transmis par SMS ou via un clavier.
Les fraudes incluent~: modification du firmware pour altérer
les mesures, utilisation de tokens volés, court-circuit du compteur.

\begin{tableaucours}{p{3.5cm}p{4.0cm}p{4.5cm}}{%
    \tableauHeader{Type de fraude} &
    \tableauHeader{Trace forensique} &
    \tableauHeader{Où la trouver}
}
Modification firmware & Firmware non signé ou corrompu &
    Extraction UART/JTAG~; vérification signature du fabricant \\
Utilisation tokens volés & Tokens non générés par SONATREL &
    Logs de la base de données tokens SONATREL (réquisition) \\
Court-circuit du compteur & Absence de données de consommation
    sur la période fraudée &
    Logs du concentrateur PLC/RF de SONATREL~; relevés voisins \\
Clonage du numéro de série & Deux compteurs avec le même SN &
    Base de données SONATREL~; inspection physique \\
\end{tableaucours}
\vspace{-0.8em}
\begin{center}{\small\itshape Tableau~18.2 --- Fraudes compteurs prépayés et traces forensiques}\end{center}

% ============================================================================
\section{Principes de Custody pour les Preuves IoT}
\label{sec:custody-iot}
% ============================================================================

\subsection{Défis spécifiques à l'IoT}

\begin{tableaucours}{p{3.5cm}p{4.0cm}p{4.5cm}}{%
    \tableauHeader{Défi} &
    \tableauHeader{Impact forensique} &
    \tableauHeader{Mitigation}
}
Absence de write blocker standardisé &
    Connexion à l'interface admin peut modifier les logs &
    Documenter chaque connexion~; ne jamais écrire via SSH \\
Volatilité extrême de la RAM &
    Les logs en RAM disparaissent en secondes &
    Exporter immédiatement~; documenter l'heure d'export \\
Horloge non synchronisée &
    Dérive de l'horloge rend les logs inutilisables &
    Mesurer la dérive NTP au moment de la saisie~;
    corriger dans la timeline \\
Firmware propriétaire &
    Difficile à analyser sans outils du fabricant &
    Binwalk~; contacter le fabricant via réquisition \\
Alimentation continue &
    Couper l'alimentation efface la RAM &
    Exporter d'abord~; couper seulement quand prêt \\
\end{tableaucours}
\vspace{-0.8em}
\begin{center}{\small\itshape Tableau~18.3 --- Défis custody IoT et mitigations}\end{center}

\subsection{Protocole de saisie IoT}

\begin{boxPNC}
\textbf{Protocole PNC --- Saisie d'un appareil IoT (routeur, caméra, DVR)}

\begin{enumerate}
    \item \textbf{Ne pas éteindre immédiatement.}
          Photographier l'état des LEDs, des câbles connectés,
          l'écran si présent. Documenter l'IP affichée.
    \item \textbf{Tenter l'accès à l'interface admin} (SSH, web,
          Telnet)~: exporter les logs volatils en priorité.
          Si l'accès échoue~: passer à l'étape~5.
    \item \textbf{Vérifier la synchronisation horaire.}
          Comparer l'heure de l'appareil avec une référence GPS.
          Documenter la dérive exacte (ex.~: ``2026-03-13~14h22
          affiché sur l'appareil~= 14h25 heure GPS, dérive~: +3~min'').
    \item \textbf{Exporter les données volatiles}~: syslog, dmesg,
          table ARP, connexions actives, baux DHCP. Calculer le hash
          SHA-256 de chaque fichier exporté.
    \item \textbf{Couper l'alimentation proprement} (via commande
          \texttt{halt} ou \texttt{poweroff} si possible)~; sinon
          débrancher l'alimentation. Ne jamais appuyer sur le
          bouton reset.
    \item \textbf{Retirer les supports amovibles} (SD Card, USB)
          \emph{avant} de couper l'alimentation~: ils constituent
          des preuves indépendantes.
    \item \textbf{Documenter et sceller~:} étiquette avec N° dossier
          N° evidence, date/heure, OPJ~; sachet antistatique
          (important pour l'électronique).
    \item \textbf{Transport}~: protection antistatique et
          antichoc~; jamais dans le même sachet que d'autres
          appareils électroniques.
\end{enumerate}
\end{boxPNC}

\separateur

\begin{boxresume}
\textbf{Ce qu'il faut retenir de ce chapitre~:}
\begin{itemize}
    \item \textbf{Ne pas éteindre en premier~:} contrairement aux PC,
          les appareils IoT stockent leurs logs principalement en RAM.
          Exporter les logs via l'interface admin \emph{avant} toute
          extinction.

    \item \textbf{Synchronisation horaire obligatoire~:} mesurer la
          dérive de l'horloge IoT par rapport à une référence GPS
          au moment de la saisie~; corriger dans la timeline. Sans
          cette correction, les logs IoT sont inexploitables pour
          la corrélation temporelle.

    \item \textbf{Binwalk~:} outil universel d'analyse de firmware.
          Identifie et extrait les systèmes de fichiers, les images
          compressées, les signatures. \texttt{binwalk -e firmware.bin}
          puis analyser le SquashFS extrait.

    \item \textbf{UART~:} accès console root sur presque tout
          appareil embarqué. Matériel~: adaptateur USB-UART ($\approx$~5~USD).
          Documenter photographiquement avant ouverture du boîtier
          (ACPO Principe~2).

    \item \textbf{DVR/NVR~:} exporter d'abord via interface web~;
          si effacement suspect~: extraire le HDD, image DD avec
          write blocker, carving vidéo avec \texttt{bulk\_extractor}
          ou \texttt{foremost}.

    \item \textbf{Vérifier la synchronisation horaire du NVR~:} une
          dérive de quelques heures invalide la corrélation avec
          les autres sources de la timeline.

    \item \textbf{Contexte camerounais~:} logs DNS des Box opérateurs
          (dnsmasq)~; historique connexions Wi-Fi (table ARP,
          \texttt{hostapd\_events.log})~; toujours compléter par
          réquisition opérateur (art.~57 Loi~2010).

    \item \textbf{Factory reset~:} n'efface pas systématiquement
          la SD Card~; n'efface pas le firmware~; laisse des traces
          dans les logs opérateur. Saisir physiquement avant
          tout reset.
\end{itemize}
\end{boxresume}

\bigskip

\begin{boxexo}[title={\ding{45}\quad Exercice~18.1 --- Priorité de collecte IoT \niveaubase}]
Vous arrivez sur la scène d'un cybercafé de Yaoundé. Un routeur TP-Link
sous OpenWRT est allumé~; l'écran du propriétaire affiche une interface
web indiquant que le routeur est connecté à Internet. Vous suspectez
qu'il a servi de point d'accès pour une fraude en ligne.
\begin{enumerate}[(a)]
    \item Listez vos actions dans l'ordre exact (6~étapes minimum),
          en justifiant chaque priorité par la volatilité des données.
    \item Quelles données exportez-vous en priorité absolue, et
          pourquoi~? Donnez les commandes exactes.
    \item Si le propriétaire tente de réinitialiser le routeur
          pendant que vous l'examinez, que faites-vous~? Quelles
          preuves subsisteraient malgré le reset~?
    \item Calculez et documentez les hachages SHA-256 des
          exports~: à quel moment dans la séquence d'actions~?
\end{enumerate}
\end{boxexo}

\begin{boxexo}[title={\ding{45}\quad Exercice~18.2 --- Analyse firmware \niveaumoyen}]
Le firmware \texttt{mtnbox\_fw\_v3.2.bin} a été extrait d'une MTN Box
saisie. \texttt{binwalk mtnbox\_fw\_v3.2.bin} révèle un SquashFS à
l'offset \texttt{0x20000}. Après extraction~:
\begin{lstlisting}[language=bash_forensique]
cat /mnt/firmware/etc/shadow
# root:$1$salt$abc123defghi:18000:0:99999:7:::
# admin:$1$salt$xyz789uvwxyz:18000:0:99999:7:::

find /mnt/firmware/ -name "authorized_keys"
# /mnt/firmware/root/.ssh/authorized_keys
cat /mnt/firmware/root/.ssh/authorized_keys
# ssh-rsa AAAAB3NzaC1yc2E... root@185.220.XX.XX
\end{lstlisting}
\begin{enumerate}[(a)]
    \item Identifiez les indicateurs de compromission dans ce firmware.
    \item Quel type d'accès la clé SSH autorisée permet-elle~?
          Quelle IP contrôle cette clé (comment le savez-vous)~?
    \item Comment vérifier si ce firmware est une version modifiée
          par rapport au firmware officiel du fabricant~? Quelle
          commande utilisez-vous~?
    \item Rédigez la constatation forensique correspondante, puis
          l'analyse avec niveau de confiance.
\end{enumerate}
\end{boxexo}

\begin{boxexo}[title={\ding{45}\quad Exercice~18.3 --- Timeline IoT multi-sources \niveauexpert}]
Dans l'affaire MBONGO (Chapitre~\ref{chap:custody}), la caméra IP
installée dans le bureau comptabilité de SONATRAV~SA est au c\oe ur
de l'investigation~: elle aurait enregistré MBONGO en train de
manipuler le système de virement. Le NVR Hikvision affiche dans
ses logs~:
\begin{lstlisting}[language=bash_forensique]
[2026-01-15 13:50:22] LOGIN  admin  192.168.5.14  Success
[2026-01-15 13:50:45] PLAYBACK  admin  Camera-Bureau-Compta
                       2026-01-15 12:30 -> 13:00
[2026-01-15 13:51:02] DELETE  admin  Camera-Bureau-Compta
                       2026-01-15 12:00 -> 14:00
[2026-01-15 13:51:08] EXPORT  admin  Camera-Couloir-1
                       2026-01-15 12:00 -> 14:00 -> USB1
\end{lstlisting}
L'heure GPS lors de la saisie du NVR est 14h28. L'heure affichée
sur le NVR est 14h21 (dérive~: --7~minutes).
\begin{enumerate}[(a)]
    \item Corrigez les horodatages NVR pour les aligner sur l'heure
          réelle. Que devient l'horodatage du DELETE~?
    \item Que révèlent les logs NVR sur les intentions de l'auteur~?
          Quelles données ont été supprimées~? Quelles données
          ont été exportées~? Pourquoi cette distinction est-elle
          significative~?
    \item Les enregistrements supprimés (12h00--14h00) sont-ils
          récupérables~? Quelle procédure appliquez-vous~?
    \item L'IP \texttt{192.168.5.14} a accédé au NVR. Comment
          identifier la machine correspondante~? Quelles sources
          corroborent cette identification~?
\end{enumerate}
\end{boxexo}

\bigskip

\begin{boxinfo}
\textbf{Clôture de la Partie~IV.} Ce chapitre clôt la Partie~IV
(Techniques Avancées). Les cinq chapitres de cette partie ont
couvert~: Forensique Système (Ch.~\ref{chap:for-sys}), Forensique
Réseau (Ch.~\ref{chap:for-res}), Forensique Mobile
(Ch.~\ref{chap:for-mob}), Forensique Cloud (Ch.~\ref{chap:for-cloud})
et ce chapitre (IoT/Embarqué).

La \textbf{Partie~V} aborde les contre-mesures~: Anti-Forensique
(Ch.~19~--- techniques d'effacement et obfuscation~; comment les
détecter et les documenter) et les investigations avancées
(Ch.~20--22). La connaissance de l'anti-forensique est indissociable
de la maîtrise des techniques forensiques présentées dans
cette Partie~IV.
\end{boxinfo}
